% met_01_flujo_cerrado — Secuencia experimental y separación de las particiones.
% Fuente editable (TikZ). Sustento: muografia-documentada, main 9558099 (protocolos EF-01 y
% OE3-C1, cierre del OE3). Compilar con: pdflatex met_01_flujo_cerrado.tex
% Diagrama sin resultados ni porcentajes de cumplimiento. Paleta Paul Tol "bright"; los
% objetivos se distinguen también en grises por el trazo del borde (continuo, discontinuo,
% punteado).
\documentclass[12pt,border=2pt]{standalone}
\usepackage[T1]{fontenc}
\usepackage[utf8]{inputenc}
\usepackage[spanish]{babel}
\usepackage{lmodern}
\usepackage{amsmath}
\usepackage{tikz}
\usetikzlibrary{babel,positioning,arrows.meta,calc}
\definecolor{tolgreen}{HTML}{228833}
\definecolor{tolblue}{HTML}{4477AA}
\definecolor{tolpurple}{HTML}{AA3377}
\definecolor{tolred}{HTML}{EE6677}
\begin{document}
\scriptsize
\hyphenpenalty=10000 \exhyphenpenalty=10000
\begin{tikzpicture}[
  >={Stealth[length=1.8mm]},
  caja/.style={draw=#1!85!black, fill=#1!10, rounded corners=1.5pt, line width=0.6pt,
               text width=24mm, minimum height=12mm, align=center, inner sep=1.6pt},
  oe1/.style={caja=tolgreen},
  oe2/.style={caja=tolblue, dashed},
  oe3/.style={caja=tolpurple, densely dotted, line width=0.85pt},
  flecha/.style={->, line width=0.6pt, draw=black!75},
  roja/.style={->, line width=0.7pt, draw=tolred!85!black},
  rot/.style={font=\tiny, text=black!75, align=center}
]
\def\dx{27.5mm}
% ---------------- Fila 1: OE1 -------------------------------------------------
\node[oe1] (cat) at (0,0) {Catálogo de 19 geometrías y fuente de muones};
\node[oe1] (g4)  at (\dx,0) {Transporte en Geant4 por bloques};
\node[oe1] (ref) at (2*\dx,0) {114 referencias: 19 geometrías por 6 vistas};
\node[oe1] (sub) at (3*\dx,0) {Submuestreo binomial: 4 fracciones y 20 realizaciones};
\node[oe1] (par) at (4*\dx,0) {Partición por geometría: 12 / 3 / 3 / 1};
\foreach \a/\b in {cat/g4, g4/ref, ref/sub, sub/par} \draw[flecha] (\a) -- (\b);

% ---------------- Fila 2: OE2 y OE3 ------------------------------------------
\def\yb{-29mm}
\node[oe2] (esc)  at (0,\yb) {Escala $s$ ajustada solo con entrenamiento};
\node[oe2] (ent)  at (\dx,\yb) {Entrenamiento al 5\,\% y selección con validación};
\node[oe2] (cong) at (2*\dx,\yb) {Ocho modelos congelados, identificados por huella};
\node[oe2] (ef)   at (3*\dx,\yb) {Evaluación EF-01 con entrada al 5\,\%};
\node[oe3, text width=27mm] (c1)   at (4*\dx,\yb) {OE3-C1: mismos ocho pesos con entrada al 1, 5, 10 y 50\,\%; bloque 5\,\% reutilizado};
\foreach \a/\b in {esc/ent, ent/cong, cong/ef, ef/c1} \draw[flecha] (\a) -- (\b);

% ---------------- Fila 3: salidas del OE3 -------------------------------------
\def\yc{-50mm}
\node[oe3] (cur)  at (3*\dx,\yc) {Curvas descriptivas de la entrada sin red en cuatro fracciones};
\node[oe3, text width=27mm] (crit) at (4*\dx,\yc) {Métricas locales y criterio conjunto por vista: SSIM y PSNR frente a la entrada sin red al 50\,\%};
\draw[flecha] (c1) -- (crit);
\draw[flecha] (ef) -- node[rot, left] {mismas entradas\\y regiones} (cur);

% ---------------- Partición: entrenamiento/validación y reservas --------------
\def\ycan{-12mm}
\draw[flecha] ($(par.south)+(-6mm,0)$) -- ($(par.south)+(-6mm,-10.5mm)$) -| (esc.north);
\node[rot, anchor=south] at ($(ent.north)+(0,7.6mm)$) {entrenamiento y validación};
\draw[roja] ($(par.south)+(6mm,0)$) -- ($(par.south)+(6mm,-6.5mm)$) -| ($(ef.north)+(5mm,0)$);
\node[rot, text=tolred!80!black, anchor=south] at ($(ef.north)+(5mm,8.6mm)$) {prueba y generalización};

% frontera: las particiones reservadas no ajustan escala ni pesos
\coordinate (fx) at ($(cong.east)!0.5!(ef.west)$);
\draw[tolred!85!black, line width=1.1pt, dash pattern=on 3.2pt off 1.6pt]
  ($(fx)+(0,9mm)$) -- ($(fx)+(0,-29mm)$);
\node[rot, text=tolred!80!black, text width=30mm, anchor=north east] at ($(fx)+(-1.5mm,-14mm)$)
  {Frontera: prueba y generalización solo se evalúan; no ajustan la escala ni los pesos};

% ---------------- Leyenda -------------------------------------------------------
\node[oe1, text width=11mm, minimum height=4.5mm, font=\tiny] (l1) at (0,\yc) {OE1};
\node[oe2, text width=11mm, minimum height=4.5mm, font=\tiny, right=2mm of l1] (l2) {OE2};
\node[oe3, text width=11mm, minimum height=4.5mm, font=\tiny, right=2mm of l2] (l3) {OE3};
\end{tikzpicture}
\end{document}
